% Lean mapping note:
% Ordinary IsSupremum s A is ambient-order supremum.  Relative notions below
% name the text-side version where upper and lower bounds are restricted to an
% ambient carrier S.

\begin{remark*}[Exposition]
The definitions of upper bound, lower bound, supremum, and infimum always
take place inside some ambient ordered universe.  In the ordinary notation
$s=\sup A$, the ambient type is usually left implicit: for a set of real
numbers, upper bounds are real upper bounds; for a set of rational numbers,
upper bounds are rational upper bounds.

Sometimes the ambient carrier matters enough to be named explicitly.  If
$A\subseteq S\subseteq T$ and $T$ is ordered, an element $u\in S$ is a
\emph{relative upper bound of $A$ in $S$} if every element of $A$ lies below
$u$.  Similarly, $\ell\in S$ is a \emph{relative lower bound of $A$ in $S$}
if $\ell$ lies below every element of $A$.
\end{remark*}

\begin{definitionbox}{Relative Bounds}
\begin{definition}[Relative Bounds]
\label{def:relative-bounds}
Let $(T,\leq)$ be an ordered set, let $A\subseteq T$, and let $S\subseteq T$.
An element $u\in S$ is a \emph{relative upper bound of $A$ in $S$} if
\[
  (\forall a\in A)(a\leq u).
\]
An element $\ell\in S$ is a \emph{relative lower bound of $A$ in $S$} if
\[
  (\forall a\in A)(\ell\leq a).
\]
\end{definition}
\end{definitionbox}

\begin{remark*}[Standard quantified statement]
\[
  u\in S\land(\forall a\in A)(a\leq u),
  \qquad
  \ell\in S\land(\forall a\in A)(\ell\leq a).
\]
\end{remark*}

\begin{remark*}[Predicate reading]
\[
  P=\mathsf{OrderedSet}(T,\leq),\qquad A\subseteq T,\qquad S\subseteq T,
\]
\[
  \mathsf{RelativeUpperBound}(u,A,S,P)
  :=u\in S\land(\forall a\in A)(a\leq u),
\]
\[
  \mathsf{RelativeLowerBound}(\ell,A,S,P)
  :=\ell\in S\land(\forall a\in A)(\ell\leq a).
\]
\end{remark*}

\begin{remark*}[Interpretation]
Relative bounds are ordinary bounds with an explicit restriction on which
candidates may serve as bounds.
\end{remark*}

\begin{dependencies}
\begin{itemize}
  \item \hyperref[def:ordered-set]{Ordered Set}
  \item \hyperref[def:subset]{Subset}
  \item \hyperref[def:real-upper-bound]{Upper Bound}
  \item \hyperref[def:real-lower-bound]{Lower Bound}
\end{itemize}
\end{dependencies}

\begin{definitionbox}{Relative Supremum and Infimum}
\begin{definition}[Relative Supremum and Infimum]
\label{def:relative-supremum-infimum}
Let $(T,\leq)$ be an ordered set, let $A\subseteq T$, and let $S\subseteq T$.
An element $s\in S$ is the \emph{relative supremum of $A$ in $S$} if $s$ is a
relative upper bound of $A$ in $S$ and
\[
  (\forall u\in S)\left[
    (\forall a\in A)(a\leq u)\Longrightarrow s\leq u
  \right].
\]
An element $i\in S$ is the \emph{relative infimum of $A$ in $S$} if $i$ is a
relative lower bound of $A$ in $S$ and
\[
  (\forall \ell\in S)\left[
    (\forall a\in A)(\ell\leq a)\Longrightarrow \ell\leq i
  \right].
\]
\end{definition}
\end{definitionbox}

\begin{remark*}[Standard quantified statement]
\[
  s\in S\land(\forall a\in A)(a\leq s)\land
  (\forall u\in S)\left[
    (\forall a\in A)(a\leq u)\Longrightarrow s\leq u
  \right],
\]
\[
  i\in S\land(\forall a\in A)(i\leq a)\land
  (\forall \ell\in S)\left[
    (\forall a\in A)(\ell\leq a)\Longrightarrow \ell\leq i
  \right].
\]
\end{remark*}

\begin{remark*}[Predicate reading]
\[
  P=\mathsf{OrderedSet}(T,\leq),\qquad A\subseteq T,\qquad S\subseteq T,
\]
\[
  \mathsf{RelativeSupremum}(s,A,S,P)
  :=
  \mathsf{RelativeUpperBound}(s,A,S,P)
  \land
  (\forall u\in S)
  \bigl(\mathsf{RelativeUpperBound}(u,A,S,P)\Rightarrow s\leq u\bigr),
\]
\[
  \mathsf{RelativeInfimum}(i,A,S,P)
  :=
  \mathsf{RelativeLowerBound}(i,A,S,P)
  \land
  (\forall \ell\in S)
  \bigl(\mathsf{RelativeLowerBound}(\ell,A,S,P)\Rightarrow \ell\leq i\bigr).
\]
\end{remark*}

\begin{remark*}[Interpretation]
Relative suprema and infima are sharp bounds after restricting the allowed
bound candidates to the ambient carrier $S$.
By default, the ambient set $S$ is assumed to contain the set under discussion:
$A\subseteq S$.  The definitions above can still be stated syntactically with
an arbitrary ambient set $S\subseteq T$ when it is useful to isolate which
candidates are allowed to serve as bounds.

The same ordered data can have different sharp bounds depending on the
ambient carrier.  For example,
\[
  A=\{q\in\mathbb{Q}:q^2<2\}
\]
has no relative supremum in $\mathbb{Q}$, but it has supremum $\sqrt{2}$ when
viewed as a subset of $\mathbb{R}$.  This is not a contradiction: the allowed
upper bounds changed from rational candidates to real candidates.
\end{remark*}

\begin{dependencies}
\begin{itemize}
  \item \hyperref[def:supremum]{Supremum}
  \item \hyperref[def:infimum]{Infimum}
  \item \hyperref[def:real-upper-bound]{Upper Bound}
  \item \hyperref[def:real-lower-bound]{Lower Bound}
\end{itemize}
\end{dependencies}
